% Part: sets-functions-relations
% Chapter: arithmetization
% Section: rationals

\documentclass[../../../include/open-logic-section]{subfiles}

\begin{document}

\olfileid{sfr}{arith}{rat}
\olsection{Van $\Int$ naar $\Rat$}
Hiervoor hebben we met wat na\"{i}eve verzamelingenleer de gehele
getallen opgebouwd uit de natuurlijke getallen. Nu bouwen we op bijna
dezelfde manier de rationale getallen op uit de gehele getallen. We
beginnen met twee dingen die we weten:
\begin{enumerate}
	\item Elk rationaal getal kan worden geschreven als $\nicefrac{i}{j}$.
	Hierbij zijn $i$ en $j$ gehele getallen en is $j$ niet nul.
	\item Dezelfde informatie uit $\nicefrac{i}{j}$ kunnen we ook
	vastleggen met het geordende paar $\tuple{ i, j}$.
\end{enumerate}
Het lijkt dus logisch om de rationale getallen op te vatten \emph{als}
geordende paren uit $\Int \times (\Int \setminus \{0_\Int\})$. Toch is
dat, net als eerder, iets te na\"ief. We willen namelijk
$\nicefrac{3}{2} = \nicefrac{6}{4}$, terwijl $\tuple{ 3, 2}\neq \tuple{ 6,
4}$. In het algemeen willen we:
\[
	\nicefrac{a}{b} = \nicefrac{c}{d} \text{ precies dan als } a \times d = b \times c
\]
Daarom leggen we voor paren uit $\Int \times (\Int \setminus
\{0_\Int\})$ vast wanneer twee van zulke paren hetzelfde voorstellen.
Zo krijgen we de volgende {equivalentierelatie}:
\[
	\tuple{ a, b }\Ratequiv \tuple{ c, d} \text{ precies dan als }a \times d = b \times c
\]
We moeten nog controleren of dit echt een equivalentierelatie is. Dat
gaat vrijwel hetzelfde als bij $\Intequiv$, en laten we als oefening over. 
\begin{prob}
Laat zien dat $\Ratequiv$ een equivalentierelatie is.
\end{prob}
Nu kunnen we de rationale getallen precies aangeven:
\begin{defn}
De rationale getallen zijn de equivalentieklassen onder $\Ratequiv$ van
paren van gehele getallen waarvan het tweede getal niet nul is. Zo'n
klasse bestaat uit alle paren die volgens deze relatie bij elkaar horen.
Dus $\Rat =
\equivclass{(\Int \times (\Int\setminus \{0_\Int\}))}{\Ratequiv}$.
\end{defn}
Net als bij de gehele getallen willen we een paar basisbewerkingen
vastleggen. Als $\equivrep{i,j}{\Ratequiv}$ de equivalentieklasse onder
$\Ratequiv$ is waarin $\tuple{i, j}$ als !!a{element} zit, definiëren we:
\begin{align*}
	\equivrep{a, b}{\Ratequiv} + \equivrep{c, d}{\Ratequiv} &= \equivrep{ad + bc,  bd}{\Ratequiv}\\
	\equivrep{a, b}{\Ratequiv} \times \equivrep{c, d}{\Ratequiv} &= \equivrep{a  c, b d}{\Ratequiv}.
\intertext{Nu leggen we vast wanneer $r \leq s$. We gebruiken dat
$r \le s$ precies dan als $s - r$ niet negatief is. Anders gezegd: $s - r$
kan worden geschreven als $\nicefrac{i}{j}$, met $i$ niet-negatief en
$j$~positief:}
	\equivrep{a, b}{\Ratequiv} \leq \equivrep{c, d}{\Ratequiv} &\text{ precies dan als }
	\equivrep{c, d}{\Ratequiv} - \equivrep{a, b}{\Ratequiv} = 
	\equivrep{i_\Int, j_\Int}{\Ratequiv}
\end{align*}
Dit geldt voor een bepaalde keuze waarbij $i \in \Nat$ en $0 \neq j \in \Nat$.
Nu moeten we controleren of deze definities echt werken zoals ze
\emph{bedoeld} zijn. Dat doen we in \olref[check]{sec}. Ze blijken
inderdaad te werken!{} Tot slot willen we een geheel getal ook \emph{als}
rationaal getal kunnen gebruiken. Daarom spreken we voor elke $i \in
\Int$ af dat $i_\Rat = \equivrep{i, 1_\Int}{\Ratequiv}$. In
\olref[check]{sec} controleren we weer dat dit allemaal klopt.
\begin{prob}
Laat zien dat $(i + j)_\Rat = i_\Rat+ j_\Rat$ en $(i \times j)_\Rat =
i_\Rat \times j_\Rat$ en $i \leq j \liff i_\Rat \leq j_\Rat$, voor alle
$i, j \in \Int$.
\end{prob}
\end{document}
